\section{Introduction}
\label{sec:introduction}

\subsection{Research Background and Problem Definition}
\label{subsec:background}

Multimodal large language models (MLLMs) generate open-ended responses from images and text, but their outputs can deviate from the input evidence or objective facts. Following related work \citep{sun2023aligning}, I adopt the following definition:
\begin{gapblock}
\textit{MLLM hallucination refers to factual content generated during perception and reasoning that contradicts the input visual content, the question, or the context, cannot be reasonably derived from them, or is inconsistent with verifiable real-world knowledge.}
\end{gapblock}

\textbf{Hallucination is not equivalent to general errors.} Calculation mistakes, incorrect formula selection, flawed logical reasoning, or option mapping errors are not labeled as hallucination unless they introduce factual content satisfying the conditions above. Conversely, even when the final answer is correct, an instance is still labeled as hallucination whenever its perception or reasoning process satisfies those conditions. I therefore define \textbf{correctness and hallucination labels}:
\[
\begin{aligned}
C_i&=
\begin{cases}
1, & \text{reference match},\\
0, & \text{reference mismatch},
\end{cases}
&\qquad
H_i&=
\begin{cases}
1, & \text{hallucination present},\\
0, & \text{hallucination absent}.
\end{cases}
\end{aligned}
\]
The labels $C$ and $H$ are independent and together define four combinations. This separation distinguishes responses that are wrong without hallucination from correct responses that still contain hallucination, and avoids using answer correctness as a proxy for hallucination.

\textbf{Uncertainty quantification (UQ)} assigns an uncertainty score to a model response using signals such as generation probability, repeated sampling, internal representations, or model self-assessment \citep{liu2025uncertainty}; the score reflects how uncertain the model is about that response. Throughout this report, a higher UQ score indicates higher uncertainty.

Uncertainty is related to hallucination, but the two are not equivalent. Uncertainty describes the stability of the model's predictive distribution or internal state, whereas hallucination describes whether the factual content of a response is supported by the input and reliable knowledge. High uncertainty can arise from image blur, ambiguous questions, or multiple plausible answers, and does not necessarily indicate hallucination; low uncertainty only means that the model has formed a stable prediction, not that the prediction is factually grounded. Prior work also shows that the relationship between UQ and hallucination varies across tasks, models, and hallucination types \citep{agnimo2026evaluating}.

A model may repeatedly produce the same hallucination with high probability because of language priors, stable but erroneous visual representations, or other systematic biases. For a given UQ method, I define this situation as \textbf{low-uncertainty hallucination (LUH)}: the response is independently judged to contain hallucination, yet the UQ method still places it in the low-uncertainty region. LUH is a false negative of hallucination detection and reflects the failure of a UQ method to recognize unsupported content that the model generates stably.

According to the uncertainty level and the presence of hallucination, model responses fall into four categories, as summarized in \Cref{fig:uncertainty_hallucination_quadrants}.

\begin{figure}[h]
\centering
\includegraphics[width=0.92\linewidth]{四种幻觉类型}
\caption{Four response categories formed by uncertainty level and hallucination status. The low-uncertainty hallucination (LUH) quadrant highlights the false-negative region that uncertainty quantification methods may fail to identify.}
\label{fig:uncertainty_hallucination_quadrants}
\end{figure}

\subsection{A Survey of Multimodal Hallucination Evaluation}
\label{subsec:hallucination_survey}

Multimodal hallucination evaluation has evolved from checking whether described objects exist to assessing whether open-ended responses are supported by evidence; the evaluation target has also expanded from the final answer to the complete response. Early studies focused mainly on object hallucination in image description. \citet{rohrbach2018object} proposed CHAIR, which compares the objects mentioned in a generated description with the objects present in the image and computes the proportion of hallucinated objects among all mentioned objects. This work pioneered an evaluation paradigm centered on visual evidence checking, but it is limited to a fixed set of object categories such as the 80 classes of MSCOCO, and its results are sensitive to instruction design and description length. With the rise of multimodal large models, \citet{li2023evaluating} proposed POPE, which converts object hallucination detection into a binary yes/no question about whether a certain object exists in the image: positive questions use objects that truly exist, while negative questions sample objects that do not appear, with random, popular, and adversarial negative sampling strategies. They found that objects appearing frequently or co-occurring often in the training data are more likely to be hallucinated, suggesting that multimodal hallucination may originate from stable language and co-occurrence priors rather than from visual perception noise alone. Comprehensive benchmarks such as MME further incorporate object existence, counting, position, and color perception subtasks into a unified polling framework \citep{fu2023mme}.

As evaluation targets became finer, object-level hallucination was summarized into three categories: category hallucination, which identifies object categories that do not exist in the image; attribute hallucination, in which the object category is correct but attributes such as color, shape, material, or quantity are described incorrectly; and relation hallucination, in which the objects and attributes are correct but the interactions or relative positions do not match the image. This taxonomy shows that hallucination is not limited to claiming the existence of nonexistent objects; it also includes incorrect descriptions of existing objects. Therefore, checking only whether the final answer matches a reference cannot cover unsupported content at the attribute or relation level.

\textbf{Open-ended hallucination evaluation.} Polling-based evaluation can hardly capture complex hallucination in open-ended responses. MMHal-Bench constructs an open-ended question set covering eight types of questions, such as object attributes, counting, and spatial relations, and uses GPT-4 as a judge to score responses on a 0 to 6 scale \citep{sun2023aligning}. This shifts hallucination detection from rule matching to having a stronger model judge whether a response is supported by evidence. HallusionBench takes a diagnostic perspective and divides questions into Visual Dependent, which can only be answered with visual context, and Visual Supplement, which can be answered without visual input while the image provides supplementary or corrective information; it systematically examines visual illusions, language priors, false premises, and image-text conflicts \citep{guan2024hallusionbench}. The Visual Supplement setting directly tests whether the model is dominated by language priors when they conflict with image evidence. In addition, GAVIE verifies the feasibility of GPT-4 for open-ended hallucination evaluation \citep{liu2023mitigating}, and HaELM attempts to train a dedicated hallucination detection model to reduce dependence on general-purpose large models \citep{wang2023evaluation}.

\textbf{Evolution of evaluation paradigms and connection to this study.} Overall, multimodal hallucination evaluation has progressed from rule-based object comparison (e.g., CHAIR), to discriminative polling (e.g., POPE and MME), and then to open-ended evaluation assisted by stronger models (e.g., MMHal-Bench, GAVIE, and HaELM); the evaluation target has expanded from whether the final answer matches to whether the complete response is supported by the input evidence. These works also show that a single accuracy rate cannot fully describe the trustworthiness of an MLLM. On the one hand, an open-ended response may include incorrect visual details even when its conclusion is correct. On the other hand, a model may answer incorrectly because of reasoning or calculation mistakes even when its visual observation is completely correct; such mistakes do not constitute hallucination under the definition in \Cref{subsec:background}. Therefore, studying whether UQ can predict hallucination requires hallucination labels that are independent of correctness, together with the visual observation, reasoning process, and final answer, so that the location of hallucination can be identified.

\subsection{From Language Model Uncertainty to Multimodal Uncertainty}
\label{subsec:from_llm_to_mllm}

\subsubsection{Uncertainty Quantification Methods for Large Language Models}

Classical machine learning typically distinguishes aleatoric from epistemic uncertainty. For open-ended language generation, recent surveys further divide uncertainty into input, reasoning, parameter, and prediction dimensions \citep{liu2025uncertainty}. Based on how the score is obtained, existing UQ methods for large language models fall into five categories.

\textbf{Generation-probability methods.} These methods use logits to compute token entropy, predictive entropy, sequence negative log-likelihood, or perplexity (PPL) \citep{malinin2021uncertainty}. They are computationally cheap, but generation probability only reflects how likely the response is under the model distribution; it does not directly indicate factual correctness.

\textbf{Model self-assessment methods.} P(True) and verbalized confidence estimate uncertainty by asking the model whether its answer is correct or by letting it directly report a confidence value \citep{kadavath2022language}. These methods work for closed-source models, but their results are sensitive to prompting and may inherit the model's knowledge bias.

\textbf{Multi-sample consistency methods.} These methods repeatedly generate responses and estimate uncertainty from answer frequency, textual similarity, or semantic consistency. Semantic Entropy clusters sampled responses by meaning and computes the cluster entropy \citep{kuhn2023semantic}. Kernel Language Entropy uses a semantic kernel to capture continuous similarity \citep{nikitin2024kernel}. Eccentricity and Degree are based on centrality measures of the response relation graph \citep{lin2023generating}, while EigenScore uses the spectral characteristics of the covariance matrix of response embeddings \citep{chen2024inside}. These methods suit open-ended responses, but they can still assign low uncertainty when the model stably repeats the same mistake.

\textbf{Internal-representation methods.} These methods use hidden states, attention, or intermediate-layer activations to compute distances or densities, or train lightweight probes to predict answer reliability. \citet{azaria2023internal} showed with linear probes that internal states contain information about statement truthfulness, but these methods usually require white-box access, and their results depend on the layer, task, and supervision labels.

\textbf{Perturbation and ensemble methods.} SPUQ perturbs the input and aggregates response changes under different conditions to estimate uncertainty \citep{gao2024spuq}. Monte Carlo dropout, Bayesian approximation, and model ensembling estimate model uncertainty by approximating the parameter posterior \citep{liu2025uncertainty}. These methods can reveal local instability, but they are computationally expensive and rely on perturbations preserving the meaning of the original task.

These five categories respectively measure generation probability, self-assessed confidence, output consistency, internal states, and intervention stability. None of them can be directly interpreted as hallucination probability without further verification.

\subsubsection{Uncertainty Quantification Methods for Multimodal Large Language Models}

The text-side UQ methods above can be applied directly to the text output of an MLLM, but they cannot easily tell whether the uncertainty originates from vision, text, multimodal alignment, or language decoding. Recent methods therefore begin to exploit visual inputs and cross-modal internal signals explicitly.

\textbf{VL-Uncertainty.} This method rewrites the question, applies Gaussian blur to the image, and computes the semantic cluster entropy of responses under the perturbed inputs \citep{zhang2024vluncertainty}. It extends Semantic Entropy from a fixed input to a neighborhood of semantic-equivalent image and text perturbations, but changing both modalities at once makes it difficult to attribute the source of the score, and blurring may remove visual information needed for the task.

\textbf{VAUQ.} This method combines prediction entropy on the original image with the entropy increase after masking high-attention visual tokens, aiming to measure whether the response depends on visual evidence \citep{park2026vauq}. It explicitly accounts for visual information, but high attention weight does not necessarily mean that a region is the evidence the response relies on, and its main evaluation target is still final-answer correctness.

\textbf{UMPIRE.} This method combines the dispersion of sampled responses in the MLLM's internal semantic space with generation probability, capturing both response diversity and multimodal consistency \citep{lau2026umpire}. Its uncertainty, however, mainly reflects variation across samples, and it treats answer correctness as the primary risk target.

\textbf{Visual Semantic Entropy.} This method perturbs only the image, keeps the text question fixed, and computes uncertainty from the weighted distance among semantic prototypes of responses \citep{huy2026visualsemanticentropy}. The study notes that overconfident visual representations make sampled responses highly consistent in meaning, causing standard Semantic Entropy to underestimate visual uncertainty; its evaluation target, however, remains answer errors and visual ambiguity.

In summary, multimodal UQ has evolved from reusing text-based scores to explicitly probing visual evidence and cross-modal representations. Whether these scores can identify hallucination under the strict definition above, however, still requires verification with independent hallucination labels.

\subsection{Common Issues in Existing Research}
\label{subsec:gaps}

Based on the above survey, I identify four interrelated issues in existing UQ research on predicting low-uncertainty hallucinations. For each issue, I describe the prevailing practice and why it matters for LUH detection.

\begin{gapblock}
\textbf{1. Mechanism assumption}\par\smallskip
\textbf{Common practice.} High generation probability or multi-sample consistency is interpreted as reliability.\par
\textbf{Why it matters for LUH.} When the same hallucination is reproduced stably, methods such as PPL and Semantic Entropy may still assign low uncertainty.
\end{gapblock}

\begin{gapblock}
\textbf{2. Evaluation target}\par\smallskip
\textbf{Common practice.} The correctness of the final answer is used as the hallucination or risk label.\par
\textbf{Why it matters for LUH.} This convention cannot distinguish wrong answers without hallucination from correct answers that contain hallucination, so the measured UQ effectiveness may only reflect error-detection performance.
\end{gapblock}

\begin{gapblock}
\textbf{3. Measurement object}\par\smallskip
\textbf{Common practice.} Only the final answer is scored, using generation probability or semantic dispersion.\par
\textbf{Why it matters for LUH.} Unsupported content introduced in the visual observation or reasoning process may not appear in the UQ score of a short final answer.
\end{gapblock}

\begin{gapblock}
\textbf{4. Perturbation validity}\par\smallskip
\textbf{Common practice.} Image blurring and text rewriting are treated as semantically equivalent perturbations.\par
\textbf{Why it matters for LUH.} These perturbations may destroy task information or introduce prompt sensitivity, causing a high score to reflect artificial information degradation rather than the hallucination risk of the original input.
\end{gapblock}

This view is also supported by a recent systematic study. \citet{agnimo2026evaluating} compared 46 UQ estimators and found that the association between uncertainty and hallucination varies significantly with the dataset, which encodes both the task and the hallucination type, and with the model; no single estimator is consistently best across datasets. That study further notes that existing work often substitutes broad correctness, robustness, or calibration targets for hallucination-specific targets, so UQ scores should not be directly regarded as general-purpose hallucination detectors. However, that study focuses on text LLMs and does not systematically examine low-uncertainty hallucinations caused by visual evidence, language priors, and multimodal fusion.

\subsection{Research Question and Experimental Pipeline}
\label{subsec:research_question}
\label{subsec:pipeline}

Based on the above survey, I formulate the following core research question:

\begin{gapblock}
\textbf{RQ: How can uncertainty quantification methods identify and predict low-uncertainty hallucination in multimodal large language models?}
\end{gapblock}

The low-uncertainty hallucination studied in this report is not a new hallucination type. It is a class of responses that existing UQ methods may miss: the response contains content unsupported by the image, question, context, or verifiable facts, while the corresponding UQ score remains low. To answer the research question, I first identify the blind spots of existing methods and then explore a more effective hallucination-specific uncertainty signal. The resulting two-stage experimental pipeline is shown in \Cref{fig:overall_research_pipeline}.

\begin{enumerate}[leftmargin=2em]
\item \textbf{Experiment~1: UQ evaluation and LUH blind-spot identification (Sections~\ref{sec:experiment1_framework}--\ref{sec:experiment1_results}).} A unified correctness-and-hallucination evaluation framework is used to examine whether existing UQ methods can detect hallucinations, identify stable low-uncertainty false-negative regions, and construct the difficult subsets needed for the next stage.

\item \textbf{Experiment~2: exploration and validation of a new uncertainty signal (Sections~\ref{sec:era_method}--\ref{sec:experiment2_results}).} Guided by the blind spots revealed in Experiment~1, this stage investigates whether the model's early dependence on visual evidence provides an additional uncertainty signal, and evaluates its effectiveness and consistency across model architectures.
\end{enumerate}

\begin{figure}[h]
    \centering
    \includegraphics[width=\linewidth]{研究路线图}
    \caption{Overall research pipeline. Experiment~1 generates and independently labels responses, evaluates three baseline UQ methods, and constructs a score-matched LUH difficult subset. Experiment~2 develops and validates the new method on this fixed subset.}
    \label{fig:overall_research_pipeline}
\end{figure}
